%%%%%%%%%%%%%%%%%%%%%%%%%%%%%%%%%%%%%%%
%% SYNTHETIC TEST FIXTURE - not a real cover letter, not a real person, not a
%% real application record. "김민준" and "가상기술 주식회사" are invented and
%% fixture@example.com is a reserved example-domain address (RFC 2606).
%% Never copy this file into cover_letters/ as the basis of a real letter.
%%
%% What it is for: cover.cls picks its fonts with hardcoded \fontspec switches
%% (cover.cls:30-31 and every command below :36), so it is the hard case for
%% Korean - every macro here re-selects Lato or Raleway by hand, and neither
%% has a Hangul glyph. This fixture puts Hangul in all of them: the name, the
%% contact line, every \lettercontent paragraph, the Raleway-wrapped bullet
%% list, the closing and the signature. Each block also keeps some Latin text,
%% because mixed-script line breaking is its own failure mode.
%%
%% Compile exactly the way a real Korean cover letter is compiled (from
%% cover_letters/, so the class's relative font paths resolve, on the same
%% engine as the English letter):
%%
%%   cd cover_letters && TEXINPUTS=../templates/korean: xelatex -interaction=nonstopmode ../tests/fixtures/korean-cover.tex
%%
%% Must come out at exactly 1 page, like every letter from this template.
%%%%%%%%%%%%%%%%%%%%%%%%%%%%%%%%%%%%%%%

\documentclass[]{cover}

% The one line that turns a Latin-only template into a Korean-capable one.
% It goes after \documentclass so ko.TeX sees the class's font setup.
\usepackage{korean-fonts}

\usepackage{fancyhdr}

\pagestyle{fancy}
\fancyhf{}

\rfoot{Page \thepage \hspace{0pt}}
\thispagestyle{empty}
\renewcommand{\headrulewidth}{0pt}
\begin{document}

%%%%%%%%%%%%%%%%%%%%%%%%%%%%%%%%%%%%%%
%     이름 / TITLE NAME
%%%%%%%%%%%%%%%%%%%%%%%%%%%%%%%%%%%%%%
\namesection{}{\Huge{김민준}}{  \href{mailto:fixture@example.com}{fixture@example.com} | +82 10 0000 0000 | 서울특별시 강남구 |  \urlstyle{same}\href{https://example.com/fixture}{링크드인}
}

%%%%%%%%%%%%%%%%%%%%%%%%%%%%%%%%%%%%%%
%     본문 / MAIN COVER LETTER CONTENT
%%%%%%%%%%%%%%%%%%%%%%%%%%%%%%%%%%%%%%

\currentdate{2026년 1월 2일}
\lettercontent{채용 담당자님께,}

\lettercontent{가상기술 주식회사의 데이터 사이언티스트 직무에 지원합니다. 이 문서는 한글 렌더링을 확인하기 위한 테스트 픽스처이며 실제 지원 서류가 아닙니다.}

\lettercontent{아래는 직무 요건과 연결되는 가상의 경험입니다. 목록은 본문과 같은 Raleway 글꼴 블록 안에 있으며, 한글은 그 안에서도 한국어 글꼴로 표시되어야 합니다.}

{\raggedright\fontspec[Path = OpenFonts/fonts/raleway/]{Raleway-Medium}\fontsize{11pt}{13pt}\selectfont
\begin{itemize}
    \item \textbf{수요 예측:} 예측 모델을 운영 환경에 배포하고 지표를 정기 보고했습니다.
    \item \textbf{파이프라인:} 일 단위 배치 작업을 재설계하고 실패 알림을 도입했습니다.
    \item \textbf{협업:} Claude Code로 코드 리뷰와 문서화를 자동화했습니다.
\end{itemize}\par}
\vspace{6pt}

\lettercontent{귀사의 데이터 조직이 공개한 기술 블로그를 읽고 지원을 결정했다는 설정으로 작성된 문단입니다. 실제 회사에 대한 사실 주장은 포함되어 있지 않습니다.}

\lettercontent{검토해 주셔서 감사합니다. 좋은 소식을 기다리겠습니다.}

\begin{flushright}
% 닫는 인사에 한국어와 영어를 함께 두어, 하나의 \fontspec 전환 안에서
% 두 문자 체계가 같이 조판되는지 확인한다.
\closing{감사합니다. Kind regards,}

\signature{김민준}
\end{flushright}
\end{document}
